%!TEX root = paper.tex

\section{$\ell_p$ Matrix Concentration Bounds}
\label{sec:ellp}

In this section we sketch proofs of the various matrix concentration
bounds that we utilize in our algorithms.

We get the core results from several earlier papers on $\ell_p$ approximation bounds.
To understand how these papers match our use of them, one should realize that they are
written in the language of approximate isomorphisms between Banach spaces.  Effectively,
they attempt to embed the column span of our matrix $\AA$ as a subspace of $L_p$.

\begin{theorem}
\label{thm:ellpsample}
Given an $n$ by $d$ matrix $\AA$
with $\ell_p$ Lewis weights $\wwbar$, for any set of sampling values $\pp_i$,
$\sum_i \pp_i = N$,
\begin{equation*}
\pp_i \geq f(d, N, p, \epsilon, \delta) \wwbar_i,
\end{equation*}
if we generate a matrix $\SS$ with $N$ rows, each chosen
independently as the $i$th standard basis vector,
times $\frac{1}{\pp_i^{1/p}}$, with probability $\frac{\pp_i}{N}$,
then with probability at least $1-\delta$ we have
\begin{equation*}
\norm{\SS \AA \xx}_1
\approx_{1 + \epsilon} \norm{\AA \xx}_1
\end{equation*}
for all vectors $\xx \in \Re^{d}$.

Valid asymptotic bounds on $d f(d, N, p, \epsilon, \delta)$ (i.e. the resulting row count,
since Lewis weights sum to $d$; here, the resulting row count itself
is plugged in as $N$) are given in the table below:

\begin{tabular}{| l | l | l |}
\hline
$p$ & $\delta$ & Sufficient row count \\ \hline
$p = 1$ & $\frac{1}{d^C}$ & $d \log(d / \epsilon) / \epsilon^2$ \\ \hline
$p = 1$ & $\frac{1}{C}$ & $d \log d / \epsilon^2$ \\ \hline
$1 < p < 2$ & $\frac{1}{C}$ & $d \log(d / \epsilon) \log(\log d / \epsilon)^2 / \epsilon^2$ \\ \hline
$p > 2$ & $\frac{1}{d^C}$ & $d^{p / 2} \log d \log(1 / \epsilon) / \epsilon^5$ \\ \hline
\end{tabular}
\end{theorem}

The last entry was proved directly in~\cite{BourgainLM89}.  For the others (the $p < 2$ cases),
these statements (that this sampling procedure is valid) are not proved directly.  Rather,
they look at the case where the Lewis weights are \emph{uniformly} small,
examining a random process choosing $\sigma_i$ independent Rademacher variables
(i.e. independently $\pm 1$ with probability $\frac{1}{2}$ each) to give
\begin{equation*}
\max_{\norm{\AA \xx}_p = 1} \left | \sum_i \sigma_i |\aa_i^T \xx_i|^p \right |.
\end{equation*}
It is worth noting that this is equivalent to taking the error of of sampling
about half the rows by unbiased coin flips.

We may, specifically, get the following statements from the existing results.

The following was shown in~\cite{Talagrand90}
\begin{lemma}[\cite{Talagrand90}]
\label{lem:talagrandL1}
There exists a constant $C$ such that if every row of $\AA$ has
Lewis weight at most $C \frac{\epsilon^2}{\log d}$, we have
\begin{equation*}
\expec{\sigma}{\max_{\norm{\AA \xx}_1 = 1}
\left | \sum_i \sigma_i \left|\aa_i^T \xx \right| \right|} \leq \epsilon.
\end{equation*}
\end{lemma}

\begin{proof}
This is implicit in \cite{Talagrand90}.  First, the proof of Proposition 1 in that paper includes a proof that this quantity is dominated by a constant times
\begin{equation*}
\expec{g}{\max_{\norm{\AA \xx}_1 = 1}
\left | \sum_i g_i ( \aa_i^T \xx ) \right |},
\end{equation*}
where the $g_i$ are independent standard Gaussian variables.  Note that this differs in two ways: the Rademacher variables are replaced with Gaussians, and
the inner absolute values have been removed (so that it is now just the max of a linear function). 
The first step is done via the comparison lemma for Radamacher processes,
which we state formally at the start of Section~\ref{sec:concentration}.
The second step is by the contraction principle, which in turn relies on the
convexity of the quantity being bounded.

This latter quantity is then bounded, under this assumption of bounded Lewis weights,
in the proof of Proposition 2.
The probability measure $\nu$, as defined on the bottom of page 366
corresponds to the probability distribution of rows proportional to
their Lewis weights.
It is then split into atoms with small Lewis weights, and shows that randomly
sampling a subset of these atoms gives the bound above.
Our setup in Lemma~\ref{lem:talagrandL1} is equivalent to this
situation after splitting.
This means our result follows from the same proof as on page 367.
This leads to a result as stated in Proposition $2$ with $n = d$,
$K(X) = O(\sqrt{\log{d}})$
and $C \frac{\epsilon^2}{\log d}$.
\end{proof}


The following was shown in~\cite{TalagrandLp}
\begin{lemma}[\cite{TalagrandLp}]
For any $p < 2$, there exists a constant $C$ such that if every row of $\AA$ has
Lewis weight at most
$C \frac{\epsilon^2}{\log (n) (\log (\log n / \epsilon))^2}$, we have
\begin{equation*}
\expec{\sigma}{\max_{\norm{\AA \xx}_p = 1}
\left | \sum_i \sigma_i \left|\aa_i^T \xx \right|^p \right|} \leq \epsilon.
\end{equation*}
\end{lemma}

\begin{proof}
This was proven in Proposition 2.3 in~\cite{TalagrandLp}.
It gives a bound of the form
\begin{equation*}
\Lambda_{F} \leq K \sqrt{ \frac{n}{M} \log{M} } \left( \log\log{M} + \log \left( \frac{M}{n} \right) \right).
\end{equation*}
where $\Lambda{F}$ was defined in Proposition 2.2 as
\begin{equation*}
\Lambda_{F} 
= \expct{\sigma} {\max_{\xx \in F_1, \norm{\xx}_p \leq 1}
\left | \sum_i \lambda_i \sigma_i \left| \xx(i) \right|^p \right|}.
\end{equation*}
Here $\lambda_i$ is $\wwbar_i / d$ where $\wwbar_i$ are the Lewis weights as defined in this paper, and $x(i)$
has been normalized (divided by $\lambda_i^{1/p}$).

Talagrand's $n$ is our $d$ and Talagrand's $M$ is our $n$.
However, Talagrand's proof is only using the $\frac{n}{M}$ as an upper bound on the Lewis weights;
thus, we can get our claim.  With a Lewis weight (by our definition) upper bound of $U$, expressions such as
$\sqrt{\frac{n \log M}{M}}$ can be replaced with $\sqrt{U \log M}$, and $\frac{1}{M \log^4 M}$ with
$\frac{U}{n \log^4 M}$.
\end{proof}
\cite{TalagrandLp} conjectures, immediately after its core claim Theorem 1.1,
that the extra $\log \log n$ factors are unnecessary for this result
(it describes them as ``truly parasitic'').  It also points out the difficulty of adapting its approach to remove them.

We will give a general (for $p \leq 2$) reduction, Lemma~\ref{lem:momentreduct}, of these moment bounds to moment bounds on the error of our described sampling procedure.

%The statement for the first entry is proven in Section~\ref{sec:concentration},
%using $l = \log d$.  The second is from \cite{Talagrand90} (note the difference
%between these two similar results: the first worsens the $\log d$ to
%$\log(d / \epsilon)$, but achieves high probability rather than constant).  The
%third comes from \cite{TalagrandLp}.  Both of these use $l = 1$: that is, they simply
%bound an expectation.

Now we may give the general reduction.  We describe a version with an extra
restriction that the required Lewis weight bound is larger than $O(1/d)$;
this is not actually necessary (and we will sketch how to avoid this).
\begin{restatable}{lemma}{momentreduct}
\label{lem:momentreduct}
There exist constants $C_1$, $C_2$ such that for any $p < 2$, if a uniform bound of $\frac{1}{g(p, n, d, \epsilon, \delta)} = \Omega(1/d)$ on
Lewis weights implies
\begin{equation*}
\expec{\sigma}{\left (\max_{\norm{\AA \xx}_p = 1} \left | \sum_i \sigma_i |\aa_i^T \xx_i|^p \right | \right )^l} \leq \epsilon^l \delta,
\end{equation*}
then sampling as described in Theorem~\ref{thm:ellpsample} with $\pp_i \geq g(p, N + C_1 d^2, d, \epsilon / C_2, \delta) \wwbar_i$ will satisfy
\begin{equation*}
\expec{\SS}{\left ( \max_{\norm{\AA \xx}_p = 1} | \norm{\SS \AA \xx}_p^p - 1 | \right )^l} \leq \epsilon^l \delta
\end{equation*}
and in particular, $\norm{\SS \AA \xx}_p \approx_{1 + \epsilon} \norm{\AA \xx}_p$ with probability at least $1-\delta$.
\end{restatable}
This is proved in Appendix~\ref{sec:reduct}.

The following was shown in~\cite{BourgainLM89}
\begin{lemma}[\cite{BourgainLM89}]
For any $p > 2$, there exists a constant $C$ such that if
if we sample $O(d^{p / 2} \log d \log(1 / \epsilon) / \epsilon^5)$ rows
with probability proportional to Lewis weights to give $\AA'$, we have
\begin{equation*}
\max_{\norm{\AA \xx}_p = 1}
\left | \norm{\AA' \xx}_p^p - 1 \right |  \leq \epsilon
\end{equation*}
with high probability.
\end{lemma}

\begin{proof}
This is implicitly shown in the proof of Theorem 7.3 in~\cite{BourgainLM89}.  In particular, line
7.27 is giving sufficient conditions for a Bernstein inequality, applied to random samples,
to imply approximation.  These samples are chosen according to the probability measure $\nu$
used in the paper, which is proportional to the Lewis weights.  We are only assuming the weights
are lower bounds for the sampling probabilities, but increasing the probabilities further can
only improve the condition of 7.27.
\end{proof}

There are several additional proofs of results similar to those we refer to here, including in Chapter 15 of
\cite{LedouxTBook} and in \cite{GuedonR07}.
